%Author1: Ganesh Ashok Satpute
%Webpage: ganeshasatpute.github.io
%Author2: Prof. Syed Abbas, IIT Mandi, India
%This file was used. https://math.mit.edu/~psh/exam/examdoc.pdf

%Author3: Jiayang Song
%This file is used to generate final exams for CS/SE110 Programming Language in MUST in Fall 2025


% For question-only, use the following command.
% \documentclass[12pt, a4paper, addpoints]{exam}

% For answers, use the following command.
\documentclass[12pt, a4paper, answers, addpoints]{exam}

\def\CourseTitle{Computer Programming}
\def\ExamName{Final Examination}
\def\CourseNumber{CS110 \& SE110 \& BE104}
\def\ExamDate{YYY/MM/DD}
\def\Duration{2}
\def\ExamTime{XX:XX - XX:XX}
\def\Term{2509}
\def\ExamDuration{2 Hours}


\usepackage[left=2cm, right=2cm, top=2cm, bottom=2cm]{geometry}

\usepackage{amsmath}
\usepackage{amsthm}
\usepackage{amsfonts}
\usepackage{amssymb}
\usepackage{mathrsfs}
\usepackage{tasks}
% \usepackage[margin=1in]{geometry}
\checkboxchar{$\Box$}
\usepackage{setspace}
%\usepackage{bbding}
%\usepackage{pifont}
\usepackage{wasysym}
\usepackage[shortlabels]{enumitem}
\renewcommand{\qedsymbol}{$\blacksquare$}
\usepackage{array} 
\usepackage{multicol}
\setlength{\columnsep}{1cm}
\usepackage{afterpage}
\usepackage{color} 
\usepackage{graphicx} % in preamble
\usepackage{listings}
\usepackage{upquote}

\lstset{
    escapeinside={(*@}{@*)},           % Allow LaTeX commands between (*@ and @*)
}

% new macros
\newcommand{\todo}[1]{\textcolor{red}{(To do): #1}}
\newcommand{\tored}[1]{\textcolor{red}{\textbf{#1}}}
\usepackage{xeCJK}
\defaultfontfeatures{Ligatures=TeX}

\setCJKmainfont[
    BoldFont={Noto Sans TC}
]{Noto Sans TC}

% Required by listings/typewriter text
\setCJKmonofont{Noto Sans TC}

\newcommand\blankpage{%
    \null
    \thispagestyle{empty}%
    \addtocounter{page}{-1}%
    \shownoanswers{\newpage}}
%% https://tex.stackexchange.com/a/203986
% Create a True-False question format
\newcommand*{\TrueFalse}[1]{%
\ifprintanswers
    \ifthenelse{\equal{#1}{T}}{%
        \textbf{TRUE}\hspace*{14pt}False
    }{
        True\hspace*{14pt}\textbf{FALSE}
    }
\else
    {True}\hspace*{20pt}False
\fi
} 
%% The following code is based on an answer by Gonzalo Medina
%% https://tex.stackexchange.com/a/13106/39194
\newlength\TFlengthA
\newlength\TFlengthB
\settowidth\TFlengthA{\hspace*{1.16in}}
\newcommand\TFQuestion[2]{%
    \setlength\TFlengthB{\linewidth}
    \addtolength\TFlengthB{-\TFlengthA}
    \parbox[t]{\TFlengthA}{\TrueFalse{#1}}\parbox[t]{\TFlengthB}{#2}}
%%

%%
\pagestyle{headandfoot}
%\extraheadheight{.25in}
\runningheader{{\CourseNumber} - {\CourseTitle}}
{}
{{\ExamName} - {\Term}} 
\firstpagefooter{}{Page \thepage\ of \numpages}{Please go on to the next page\ldots} %
\runningfooter{}{Page \thepage\ of \numpages}{\iflastpage{End of exam.}{Please go on to the next page\ldots}}

\renewcommand{\fillinlinelength}{1cm}  % change default length

\renewenvironment{solution}
  {\par\noindent\textbf{Solution: }\color{red}} % begin code
  {\par}    

% --- show solution chunks only when 'answers' is active ---
\usepackage{xcolor} % for \color

% Inline macro: use as \showsol{<solution text>}
\newcommand{\showsol}[1]{\ifprintanswers{\textbf{\color{red}\underline{Solution:} }\color{red}#1}\fi}

\newenvironment{show_code_sol}
  {\par\noindent\textbf{Solution:}\par
   \lstset{basicstyle=\ttfamily\small\color{red}}}
  {\par}


\usepackage{listings}
\usepackage{comment} % for \excludecomment (safe with verbatim)
% \usepackage[T1]{fontenc}


% Environment version that mimics 'solution' but guarded
\newenvironment{solonly}
  {\ifprintanswers\par\noindent\textbf{Solution: }\color{red}}
  {\par\fi}

% If you want to keep using the built-in 'solution' name,
% make sure it stays conditional:
\renewenvironment{solution}
  {\ifprintanswers\par\noindent\textbf{Solution: }\color{red}}
  {\par\fi}

\newcommand{\shownoanswers}[1]{\ifprintanswers\else #1\fi}

% Base style for C
% One listings + comment load
\usepackage{listings}
\usepackage{comment} % safe with verbatim

% --- styles (no frame, no line numbers) ---
\lstdefinestyle{examc}{
  language=C,
  basicstyle=\ttfamily\small,
  numbers=none,
  frame=none,
  columns=fullflexible,
  keepspaces=true,
  showstringspaces=false,
  tabsize=4,
  breaklines=true,
  keywordstyle=\bfseries,
  commentstyle=\itshape,
  stringstyle=\ttfamily
}
\lstdefinestyle{examc-sol}{
  style=examc,
  basicstyle=\ttfamily\small\color{red}
}

% --- Environments ---
\lstnewenvironment{Ccode}[1][]
  {\lstset{style=examc,#1}}
  {}

% Define CcodeSol only when answers are printed; otherwise exclude it.
\ifprintanswers
  \lstnewenvironment{CcodeSol}[1][]
    {\lstset{style=examc-sol,#1}}
    {}
\else
  \excludecomment{CcodeSol} % body is ignored; no verbatim processing
\fi

% Inline C
\newcommand{\Cinline}[1]{\lstinline[language=C]{#1}}

% ---------------------------------------------------------------------
  


%%%%%%%%%%%%%%%%%%%%%%%%%%%%
\begin{document}

\noindent
\begin{minipage}[t]{0.48\textwidth}
    \vspace{-10pt}
    \fbox{\parbox[c][1.2cm][c]{4cm}{Seat\\No.}}
\end{minipage}
\hfill
\begin{minipage}[t]{0.48\textwidth}
    \raggedleft
    \vspace{-30pt}
    \includegraphics[width=0.35\textwidth]{must-logo.pdf}
\end{minipage}


\begin{center}
    
    \vspace*{0.5cm}
    % \includegraphics[width=0.2\textwidth]{must-logo.pdf}\\[1em]
    
    {\textbf{\Huge 澳\hspace{2em}門\hspace{2em}科\hspace{2em}技\hspace{2em}大\hspace{2em}學}}\\[0.3cm]
    {\large \textbf{MACAU UNIVERSITY OF SCIENCE AND TECHNOLOGY}}\\[1cm]

    {\Large \textbf{School of Computer Science and Engineering}}\\[0.2cm]
    {\Large \textbf{Faculty of Innovation Engineering}}\\[0.2cm]
    {\Large \textbf{Bachelor of Science}}\\[1cm]

    {\large \textbf{2025/2026 – 1st Semester {\color{red} Make Up} Exam Question Booklet}}\\[1cm]

    {\Large \textbf{\CourseNumber \quad \CourseTitle}}\\[1cm]
\end{center}

% \noindent
% \hspace{2cm}% adjust indent width here
\begin{center}
\begin{tabular}{@{}p{4cm}p{8cm}@{}}
    \textbf{No. of Pages:} & {\numpages} Pages (including this page)\\[0.5cm]
    \textbf{Exam Date:} & \ExamDate  \\[0.5cm]
    \textbf{Exam Duration:} & {\ExamDuration} ({\ExamTime}) \\[0.5cm]
    \textbf{Student Name:} & \underline{\hspace{7cm}}\\[0.5cm]
    \textbf{Student Number:} & \underline{\hspace{7cm}}\\
    % Seat No.: & \underline{\hspace{7cm}}\\
\end{tabular}
\end{center}

\vspace{0.5cm}

\noindent
\textbf{Please read the following instructions carefully.}
\begin{itemize}[leftmargin=*]
    \item This question booklet contains \textbf{\numquestions} questions, \textbf{\numpages} pages (including the cover) for the total of \textbf{\numpoints} points/marks. Check whether all pages are present.
    
    \item Please write your answers on the Question Booklet.
    
    \item Return this Question Booklet, together with all extra unused answer sheets and draft paper, at the end of the examination.
    
    \item This is a \textbf{Closed Book Exam}. The use of dictionaries, calculators, or electronic devices is \textbf{NOT permitted}.
    
    \item \textbf{Answers must be provided in English to receive marks; failure to do so will result in no marks being awarded.}

    \item \textbf{Read the instructions for individual questions carefully} before answering the questions. Clearly identify your final answers.
\end{itemize}


\vfill

\hspace{0.5in}
\begin{questions}
\addpoints


\newpage

% Q1 (2)
\question[2]{In C, parameters are passed by value. Explain what this means, and why the statement \texttt{swap(a, b);} in the function below fails to swap two integers in \texttt{main()}.
\begin{lstlisting}
    void swap(int x, int y){ 
        int t=x; 
        x=y; 
        y=t; 
    }
\end{lstlisting}
}

\shownoanswers{\vspace{4cm}}

\showsol{
C uses pass-by-value: a function receives \emph{copies} of arguments. \texttt{swap} modifies only its local copies \texttt{x} and \texttt{y}, leaving \texttt{a} and \texttt{b} in \texttt{main()} unchanged. To swap, pass \emph{addresses} and dereference.

\textbf{Remark: Full marks if “copies” is explicitly stated.}
\vspace{1cm}
}

% Q2 (2)
\question[2]{Given \texttt{int v[6] = \{5, 15, 25, 35, 45, 55\}; int *p = v + 1;} what does \texttt{*(p + 2)} evaluate to? Briefly explain the pointer arithmetic.}

\shownoanswers{\vspace{4cm}}

\showsol{
\textbf{35}

\texttt{p} points to \texttt{v[1]}. \texttt{p+2} points to \texttt{v[3]}, but \texttt{*(p+2)} dereferences element two positions after \texttt{v[1]}, i.e., \texttt{v[3]} which is 35. \\

\textbf{Remark: 1pt correct value, 1pt correct explanation. }
\vspace{1cm}
}

% Q3 (2)
\question[2]{Explain the difference between the character constant \texttt{'\textbackslash0'} and the string literal \texttt{""} (empty string).}

\shownoanswers{\vspace{4cm}}

\showsol{
\texttt{'\textbackslash0'} is a single character constant (type \texttt{int} in C), value 0. \texttt{""} is a string literal: an array of one \texttt{char} containing only \texttt{'\textbackslash0'}.

\textbf{Remark: Give full marks if array vs. single character is noted.}
\vspace{1cm}
}

% Q4 (2)
\question[2]{Is there any potential bugs in the following code? Justify your answer.
\vspace{-20pt}
\begin{center}
    \begin{lstlisting}[basicstyle=\small\ttfamily]
    #include <stdio.h>
    int main(void) {
        char c;
        while ((c = getchar()) != EOF)
            putchar(c);
        return 0;
    }
    \end{lstlisting}
\end{center}
}

\shownoanswers{\vspace{3.5cm}}

\showsol{
Potential bug: \texttt{c} is declared as \texttt{char}, but \texttt{EOF} is typically \texttt{-1} (int). 
Assigning \texttt{EOF} to \texttt{char} can truncate and break the loop condition. \\
\textbf{Fix:} use \texttt{int c;} instead of \texttt{char c;} so \texttt{EOF} is representable.

\textbf{Remark: Give full marks if the student mentions that \texttt{EOF} is typically \texttt{-1} (an \texttt{int}), or that \texttt{EOF} cannot be represented by a \texttt{char}. Reflect CO4.}
\vspace{1cm}
}

% Q5 (2)
\question[2]{Explain the difference between \texttt{break} and \texttt{continue} in loop control.}

\shownoanswers{\vspace{4cm}}

\showsol{
\texttt{break} exits the nearest loop immediately; \texttt{continue} skips to the next iteration. 

\textbf{Remark: must state effects on control flow. Reflect CO1.}
}

% Q6 (2)
\question[2]{What does this declaration mean? \texttt{char (*pa)[5];}}

\shownoanswers{\vspace{4cm}}

\showsol{
\texttt{pa} is a \emph{pointer to an array of 5 \texttt{char}}.

\textbf{Remark: 1pt “pointer to array”, 1pt array size/type.}
\vspace{1cm}
}

% Q7 (2)
\question[2]{Explain the difference between a \texttt{struct} and a \texttt{union}.}

\shownoanswers{\vspace{4cm}}

\showsol{
A \texttt{struct} stores all members simultaneously (total size \(\ge\) sum with padding). 
A \texttt{union} reuses the same memory for all members (size is the max member), only one valid at a time. 

\textbf{Remark: memory layout difference must be stated.}
\vspace{1cm}
}

% Q8 (2)
\question[2]{What happens if you open a file with mode \texttt{"a"} using \texttt{fopen()} and the file already exists?}

\shownoanswers{\vspace{4cm}}

\showsol{
Mode \texttt{"a"} opens for writing at the end; the file is \emph{not} truncated, and writes are appended to existing contents (created if missing).

\textbf{Remark: Full mark for "appends in the end".}
\vspace{1cm}
}

% Q9 (3)
\question[3]{In function parameters, \texttt{int a[]} and \texttt{int *a} appear interchangeable. Are they the same? State one similarity and one difference. }

\shownoanswers{\vspace{4cm}}

\showsol{
No, they are close but not the same.\\
Similarity: In parameter lists, \texttt{int a[]} is adjusted to \texttt{int *a}; both behave as a pointer to the first element. \\
Difference: Only the first dimension decays; information about the array's length is lost. Also, outside the parameter context, arrays are not pointers (e.g., \texttt{sizeof} differs). 
\vspace{1cm}
}

% Q10 (3)
\shownoanswers{\newpage}
\question[3]{Consider the program below:
\vspace{-20pt}
\begin{center}
\begin{lstlisting}[basicstyle=\small\ttfamily]
#include <stdio.h>
int main(void){
    unsigned int u = 4294967290U; // assume 32-bit unsigned
    int s = (int)u;
    printf("%d\n", s);
    return 0;
}
\end{lstlisting}
\end{center}
Explain why the printed value is negative on typical 32-bit two's complement systems.}

\shownoanswers{\vspace{4cm}}

\showsol{
The cast from large \texttt{unsigned} to \texttt{signed int} reinterprets the same bit pattern; since the top bit is 1, it represents a negative value in two’s complement, yielding a negative result (implementation-defined but commonly negative).

\textbf{Remark: mention sign bit and integer overflow. Reflect CO2.}
\vspace{1cm}
}

% Q11 (3)
\question[3]{What will the following code print? (Show the full output.)
\vspace{-20pt}
\begin{center}
\begin{lstlisting}[basicstyle=\small\ttfamily]
#include <stdio.h>
int main(void) {
    int i = 1;
    while (i <= 6) {
        printf("%d ", i++);
        i++;
    }
    return 0;
}
\end{lstlisting}
\end{center}
}

\shownoanswers{\vspace{3cm}}

\showsol{
\textbf{1 3 5 } \\
Each iteration prints \texttt{i} then \texttt{i++}; the extra \texttt{i++} in the body skips every other number until the condition fails at 7.

\textbf{Remark: 1pt per correct number. Reflect CO3.}
\vspace{1cm}
}

% Q12 (3)
\question[3]{Consider the program below:
\vspace{-20pt}
\begin{center}
\begin{lstlisting}[basicstyle=\small\ttfamily]
#include <stdio.h>

void set_zero(int *p){ *p = 0; }

int main(void){
    int x = 42;
    set_zero(x);
    printf("%d\n", x);
    return 0;
}
\end{lstlisting}
\end{center}
What will this program output? If it does not work as intended, how can it be fixed?}

\shownoanswers{\vspace{4cm}}

\showsol{
A Compilation error will occur: \texttt{set\_zero} expects a pointer; \texttt{x} is passed by value, causing a type mismatch. \\
\textbf{Fix:} call \texttt{set\_zero(\&x);} so the function can modify \texttt{x} via its address.

\textbf{Remark: 1pt identifying issue, 2pts correct fix.}
\vspace{1.5cm}
}

% Q13 (6) parts
\question[4]{Consider the program below:
\vspace{-20pt}
\begin{center}
\begin{lstlisting}[basicstyle=\small\ttfamily]
#include <stdio.h>
int main(void){
    char s[20] = "Hello World";
    char *p = s;
    while (*p){
        if (*p == 'o') *p = 'a';
        p++;
    }
    printf("%s\n", s);
    return 0;
}
\end{lstlisting}
\end{center}
\begin{parts}
\part(2 points) What is printed?
\shownoanswers{\vspace{3cm}}

\showsol{\texttt{Hella Warld}\\
\textbf{Remark: Reflect CO5.}
\vspace{0.5cm}}

\part(2 point) How many characters are in \texttt{s} (excluding \texttt{'\textbackslash0'}) after modification?

\shownoanswers{\vspace{3cm}}
\showsol{11 characters.\\
\textbf{Remark: Reflect CO4.}
\vspace{0.5cm}}

% \part(3 points) If we change to \texttt{char *s = "Hello World";}, what runtime issue may occur and why?
% \shownoanswers{\vspace{3cm}}

% \showsol{String literals may be stored in read-only memory; writing through \texttt{s} is undefined behaviour and may crash.\vspace{0.5cm}}

\end{parts}
}

% Q14 (4)
\shownoanswers{\newpage}
\question[4]{Write the exact output (including spaces/newlines):
\vspace{-20pt}
\begin{center}
\begin{lstlisting}[basicstyle=\small\ttfamily]
#include <stdio.h>
void rev(const char *s){
    if (*s == '\0') return;
    rev(s + 1);
    putchar(*s);
}
int main(void){
    const char *t = "go go";
    rev(t);
    putchar('\n');
    return 0;
}
\end{lstlisting}
\end{center}
}

\shownoanswers{\vspace{3cm}}

\showsol{
\texttt{og og}\\
\textbf{Remark: reversed string, space preserved.}
\vspace{1cm}
}

% Q15 (4)
\question[4]{Write the exact output:
\vspace{-20pt}
\begin{center}
\begin{lstlisting}[basicstyle=\small\ttfamily]
#include <stdio.h>
int main(void){
    int a[4] = {2, 4, 6, 8};
    int *pp[2] = { &a[0], &a[2] };
    int (*pa)[4] = &a;
    printf("%d %d %d %d\n", *pp[0], *pp[1], (*pa)[3], *(*pa + 1));
    return 0;
}
\end{lstlisting}
\end{center}
}

\shownoanswers{\vspace{3cm}}

\showsol{
\texttt{2 6 8 4}\\
\textbf{Remark: 1pt each correct value.}
\vspace{1.5cm}
}

% Q16 (5) fill blanks
\shownoanswers{\newpage}
\question[5]{Fill the blanks to read a word and compute its length using a pointer-only \texttt{mystrlen()}.
\textbf{Example Input:} \texttt{Kernel} \\
\textbf{Example Output:} \texttt{Length = 6}
\vspace{-10pt}
\begin{center}
\begin{lstlisting}[basicstyle=\small\ttfamily]
#include <stdio.h>

int mystrlen(char *s){
    int len = 0;
    char *p = (*@\fillin[\tored{s}]@*);             \\ Fill blank
    while ((*@\fillin[\tored{*p}]@*) != '\0') {     \\ Fill blank
        len++;
        (*@\fillin[\tored{p++}]@*);                 \\ Fill blank
    }
    return len;
}

int main(void){
    char buf[100];
    scanf("%99s", (*@\fillin[\tored{buf}]@*));    \\ Fill blank
    printf("Length = %d\n", mystrlen((*@\fillin[\tored{buf}]@*)));  \\ Fill blank
    return 0;
}
\end{lstlisting}
\end{center}
}
\shownoanswers{\vspace{2cm}}

\showsol{
\textbf{Remark: 1 point per correct blank.}
\vspace{1cm}
}

% Q17 (6) file sum fill blanks
\question[6]{Complete the blanks to read integers from \texttt{data.txt} and print their sum.

Assume \texttt{data.txt} contains \texttt{10  -3  7  0}\\
\textbf{Expected Output:} \texttt{Sum = 14}
\vspace{-10pt}
\begin{center}
\begin{lstlisting}[basicstyle=\small\ttfamily]
#include <stdio.h>
int main(void){
    FILE *fp = fopen((*@\fillin[\tored{"data.txt"}][2.5cm]@*), (*@\fillin[\tored{"r"}]@*));        \\ Fill blank
    if (!fp){ printf("Cannot open file.\n"); return 1; }    
    int x, sum = 0;
    while (fscanf((*@\fillin[\tored{fp}]@*), (*@\fillin[\tored{"\%d"}]@*), (*@\fillin[\tored{\&x}]@*)) == 1){       \\ Fill blanks
        sum += x;
    }
    fclose((*@\fillin[\tored{fp}]@*));      \\ Fill blank
    printf("Sum = %d\n", sum);
    return 0;
}
\end{lstlisting}
\end{center}
}
\shownoanswers{\vspace{2cm}}

\showsol{
\textbf{Remark: 1 point per correct blank.}
\vspace{1cm}
}

% Q18 (4) typedef struct pointer
\shownoanswers{\newpage}
\question[4]{Fill the blanks so the program defines a \texttt{Book} type and accesses its members via a pointer.

\textbf{Expected Output:} \texttt{Title: C Primer, Price: 59}
\vspace{-10pt}
\begin{center}
\begin{lstlisting}[basicstyle=\small\ttfamily]
#include <stdio.h>

typedef struct {
    char title[32];
    int price;
} (*@\fillin[\tored{Book}]@*);      \\ Fill blank

int main(void){
    Book b1 = {"C Primer", 59};
    Book *bp = (*@\fillin[\tored{\&b1}]@*);     \\ Fill blank
    printf("Title: %s, Price: %d\n", (*@\fillin[\tored{bp->title}][2cm]@*), (*@\fillin[\tored{bp->price}][2cm]@*)); \\ Fill blanks
    return 0;
}
\end{lstlisting}
\end{center}
}
\shownoanswers{\vspace{2cm}}

\showsol{
\textbf{Remark: 1 point per correct blank.}
\vspace{1cm}
}

% Q19 (8) transpose different matrix
\shownoanswers{\newpage}
\question[10]{Write a C program to print the transpose of the given 2D array.
The transpose of a matrix is a new matrix obtained by swapping its rows and columns. For example:

The transpose of a matrix
 $\begin{bmatrix}
  1 & 2 & 3\\ 
  4 & 5 & 6\\
\end{bmatrix}$ is 
 $\begin{bmatrix}
  1 & 4\\ 
  2 & 5\\
  3 & 6
\end{bmatrix}$.
The transpose of a matrix
 $\begin{bmatrix}
  1 & 2 \\ 
  3 & 4 \\
\end{bmatrix}$ is 
 $\begin{bmatrix}
  1 & 3\\ 
  2 & 4\\
\end{bmatrix}$.

\textbf{Expected Output:}
\(
\begin{matrix}
1 & 4 & 7\\
0 & 5 & 8
\end{matrix}
\)

\textbf{Program to Complete:}
\begin{lstlisting}[basicstyle=\small\ttfamily]
#include <stdio.h>
int main(void){
    int m[3][2] = {{1,0},{4,5},{7,8}};
    int r = 3, c = 2;
    // Print transpose (2 x 3)
    // Your code here
}
\end{lstlisting}
}

\begin{CcodeSol}
#include <stdio.h>
int main(void){
    int m[3][2] = {{1,0},{4,5},{7,8}};
    int r = 3, c = 2;
    for (int j = 0; j < c; ++j){
        for (int i = 0; i < r; ++i){
            printf("%d ", m[i][j]);
        }
        putchar('\n');
    }
    return 0;
}
\end{CcodeSol}

\showsol{
\textbf{Marking guidance (suggested):}
\begin{itemize}
\item Up to 4 pts deduction for syntax errors.
\item 2 pts for correct nested loops.
\item 1 pt for correct element access \texttt{m[i][j]}.
\item 1 pt for printing newlines per row.
\end{itemize}
\vspace{1cm}
}

% Q20 (10) sort array pointer arithmetic only
\shownoanswers{\newpage}
\question[10]{Complete function \texttt{sort}. Print in descending order.

\textbf{Expected Output:} \texttt{The sorted array is: 11 2 0 -3 -10  }

\textbf{Program to Complete:}

\begin{lstlisting}[basicstyle=\small\ttfamily]
#include <stdio.h>
void sort(int *, int);

int main(void){
    int a[] = {2, -3, 11, 0, -10};
    int N = sizeof(a)/sizeof(a[0]);
    sort(a, N);
    printf("The sorted array is: ");
    for (int i = 0; i < N; ++i) printf("%d ", a[i]);
    putchar('\n');
    return 0;
}

void sort(int *p, int n){
    // Complete function
}
\end{lstlisting}
}

\begin{CcodeSol}
#include <stdio.h>
void sort(int *, int);

int main(void){
    int a[] = {2, -3, 11, 0, -10};
    int N = sizeof(a)/sizeof(a[0]);
    sort(a, N);
    printf("The sorted array is: ");
    for (int i = 0; i < N; ++i) printf("%d ", a[i]);
    putchar('\n');
    return 0;
}

void sort(int *p, int n){
    // insertion sort using pointers
    for (int i = 1; i < n; ++i){
        int key = *(p + i);
        int j = i - 1;
        while (j >= 0 && *(p + j) <= key){
            *(p + j + 1) = *(p + j);
            --j;
        }
        *(p + j + 1) = key;
    }
}
\end{CcodeSol}

\showsol{
\textbf{Suggested marking:}
\begin{itemize}
\item Deduct up to 4 pts for syntax errors.
\item 2 pts for loop structure.
\item 2 pts for correct ordering (ascending).
\end{itemize}
\vspace{1cm}
}

% Q21 (10) struct
\question[10] {Write a C program that defines a structure called \texttt{Student} to store each student's name and score.  
You must write a function that computes and returns the \textbf{highest score} in the array of structures, and prints the \textbf{name and score} of the top student.
The initial code is given; complete this program.

\textbf{Expected Output:}

\texttt{Top Student: Carol (95)}


\textbf{Program to Complete:}

\begin{lstlisting}[basicstyle=\small\ttfamily]
#include <stdio.h>
#include <string.h>

struct Student {
    char name[50];
    int score;
};

int find_top(const struct Student *, int);

int main(void) {
    struct Student s[] = {
        {"Carol", 95},
        {"Alice", 90},
        {"Bob", 85}
    };
    int N = sizeof(s) / sizeof(s[0]);

    int index = find_top(s, N);
    printf("Top Student: %s (%d)\n", s[index].name, s[index].score);

    return 0;
}

int find_top(const struct Student *p, int n) {
    // Complete this function
}
\end{lstlisting}

}


\begin{CcodeSol}

#include <stdio.h>
#include <string.h>

struct Student {
    char name[50];
    int score;
};

int find_top(const struct Student *, int);

int main(void) {
    struct Student s[] = {
        {"Carol", 95},
        {"Alice", 90},
        {"Bob", 85}
    };
    int N = sizeof(s) / sizeof(s[0]);

    int index = find_top(s, N);
    printf("Top Student: %s (%d)\n", s[index].name, s[index].score);

    return 0;
}
int find_top(const struct Student *p, int n) {
    int max_index = 0;              // assume first student is top
    for (int i = 1; i < n; i++) {
        if (p[i].score > p[max_index].score) {
            max_index = i;
        }
    }
    return max_index;
}

\end{CcodeSol}

\showsol{    
    The following marking guidance is for reference only; you can make your own judgment. Reflect CO2.
    \begin{itemize}
        \item If students have syntax errors, deduct up to 4 points.
        \item Give 2 points if the students show the idea of using nested \texttt{for loops}, even if the loop may have mistakes.
        \item Give 2 points if the students know the idea of keeping track of the index with the highest score, and return that index to \texttt{main()}.
    \end{itemize}
    \vspace{1cm}
}

\end{questions}

\vfill
\noindent\makebox[2.5in]{\textbf{\Large Score:} \hrulefill}

\end{document}


